% META: {"id":"1SPE-SUITES-CO-051","chapitre":"1SPE-SUITES","type_objet":"corrige","exercice_ref":"1SPE-SUITES-EX-051","capacites_codes":["C8"],"sources_inspiration":[],"status":"generated"}
% BEGIN-VERIFY
% from fractions import Fraction
% assert 1 + Fraction(1, 11) == Fraction(12, 11)
% assert 3*10 + 2 == 32
% assert 2 + (-1)**10 == 3
% END-VERIFY
\begin{corrige}{1SPE-SUITES-EX-051}
\textbf{1.} Termes aux rangs demandés :

\begin{center}
\begin{tabular}{c|ccccc}
$n$ & $0$ & $1$ & $2$ & $5$ & $10$ \\ \hline
$a_n$ & $2$ & $\frac32$ & $\frac43$ & $\frac76$ & $\frac{12}{11}$ \\
$b_n$ & $2$ & $5$ & $8$ & $17$ & $32$ \\
$c_n$ & $3$ & $1$ & $3$ & $1$ & $3$
\end{tabular}
\end{center}

\textbf{2.} Ce que montre chaque représentation : les points $(n;a_n)$ ont des
ordonnées décroissantes qui se resserrent ; les points $(n;b_n)$ sont alignés et
leurs ordonnées augmentent de $3$ à chaque rang ; les points $(n;c_n)$ alternent
entre les droites d'ordonnées $3$ et $1$.

\textbf{3.} On associe donc : $(a_n)$ semble se rapprocher d'une valeur finie,
et l'on conjecture la valeur $1$ ; $(b_n)$ croît sans borne apparente ; $(c_n)$
ne se rapproche d'aucune valeur unique.

\textbf{4.} Ces conclusions reposent ici sur des observations numériques et
graphiques, \textbf{sans démonstration formelle}.
\end{corrige}
